\documentclass[10pt]{whitepaper}
\usepackage{url,graphicx,multirow,hyperref}
\hypersetup{
  bookmarks=true,
  unicode=true,
  pdftoolbar=true,
  pdfmenubar=true,
  pdffitwindow=true,
  pdfstartview={FitV},
  pdftitle={},
  pdfauthor={},
  pdfnewwindow=true,
  colorlinks=false,
  pdfdisplaydoctitle=true,
  pdfborder={0 0 0},
  pdfinfo={
    Title={Enabling MOOC Collaborations Through Modularity},
    Author={Geoffrey Challen},
  },
}
\usepackage[all]{hypcap}

\input{.xxxnote}
% 16 Nov 2010 : GWA : Any special macros or other stuff for this particular
%               whitepaper go here.
\newcommand{\opsclass}{\href{http://www.ops-class.org}{\texttt{ops-class.org}}}

\pagestyle{plain}


\def\thetitle{Enabling MOOC Collaborations Through Modularity}
\def\shorttitle{Modular MOOCs}
\def\theauthors{Geoffrey Challen (\textit{University at Buffalo}) and Margo
Seltzer (\textit{Harvard University})}
\def\shortauthors{Challen}
\def\submissiondate{04/07/2014}
\def\whitepapername{Learning with MOOCs}

\begin{document}

\ifdefined\isblue
\begin{textblock}{1}(5.9,-0.3)
\noindent\href{http://blue.cse.buffalo.edu}{\includegraphics[width=0.6in]{./figures/logos/blue.jpg}}
\end{textblock}
\else
\fi

\pagestyle{whitepaper}
\thispagestyle{emptywhitepaper}
\pagenumbering{arabic}

\chapter{\thetitle}

Massive open online courses (MOOCs) have the potential to transform teaching
and learning, and their impact is already being felt on university campuses
and in popular culture. To meet the diverse learning needs of students at
many institutions, MOOCs must encourage collaboration, not confrontation,
with the professors using them. Today's online courses, typically created and
controlled by small groups of faculty, both fail to foster this collaboration
and also miss the opportunity to leverage the experience and experimental
capacity present in classrooms across the country. A variety of instructors
teaching a variety of different students in a variety of different ways
represents an incredible education resource and living laboratory that MOOCs
must find a way to harness, not suppress.


Unfortunately, dividing lines are already being drawn
between faculty able to create MOOCs and those encouraged or even forced to
use them. In their well-publicized rejection of Harvard's popular EdX course
on social justice, philosophy faculty at San Jose State University raised
concerns over the ability of online content to integrate into their existing
courses:

\begin{quote}

When a university such as ours purchases a course from an outside vendor, the
faculty cannot control the design or content of the course; therefore we
cannot develop and teach content that fits with our overall curriculum and is
based on both our own highly developed and continuously renewed competence
and our direct experience of our students' needs and abilities.

\end{quote}

As computer scientists, we see this complaint as expressing not a fundamental
limitation of MOOCs, but a design flaw in how they are currently built:
\textbf{today's MOOCs are not modular}. Modularity has long been appreciated
when designing and implementing computer systems. Monolithic ``blobs'' of
code, usually produced by a single developer to quickly solve a particular
problem, are difficult to improve, test, and maintain, and are unlikely to be
used by anyone but their creator. All of today's major online education
providers, including Coursera, Udacity, and the edX platform, provide
monolithic MOOC content: entire courses that are packaged and offered as an
indivisible unit. This flaw limits their ability to partner with faculty at
other institutions, to generate large and diverse data sets, and to fully
address diverse student learning styles and objectives.

Since today's MOOCs suffer from a common software design problem, we are
applying a standard software engineering solution: modularization. Modular
MOOCs break content into small \textit{modules} all sharing a common
structure. Individual modules can be tested and improved in isolation and
their impact on the entire class determined. As the goals and expectations of
the class for which they were originally created change, and as the content
itself changes, it is more likely that the functionality provided by modules
will continue to remain relevant, while adding new modules keeps the course
current.

Most importantly, modularization encourages the interinstitution
collaboration that today's monolithic MOOCs struggle to achieve. Modules
allow faculty to retain full control over course content and structure, with
the ability to add, remove, alter and reorder modules in ways expressing
their understanding of their students' needs and abilities, all while
producing experimental data that benefits the original course. When faculty
provide more effective replacements for existing modules, everyone benefits.
And if one faculty member discovers an effective sequence of modules, all
participating faculty can decide whether to adjust their own courses to
incorporate this insight.

We are currently exploring the design and construction of \textit{modular
MOOCs (mMOOCs)}, which exploit modularity to provide the control traditional
faculty desire while providing the diverse and well-structure data sets
required to answer fundamental questions about teaching and learning. We are
beginning by augmenting \opsclass{}, an online operating systems
instructional environment. Built by Geoffrey Challen, it has made automated
version of the programming assignments developed by Margo Seltzer available
online. Once \opsclass{} is able to support the diverse student outcomes and
needs at our two institutions, the University at Buffalo and Harvard
University, we will recruit faculty from other institutions to foster vibrant
multi-institutional collaborations in education, just as we currently do in
research.

\vspace*{-1in}

\end{document}
